\addcontentsline{toc}{subsubsection}{Solution 5. The Triadic Balance}

\begin{dlsummary}[title={\bfseries Solution 5. The Triadic Balance},breakable]

\begin{proof}
\textbf{Step 1. Unlock the value at $3$.}

Taking $x=y=1$ in submultiplicativity gives
\[
f(1)\le f(1)^2,
\]
so $f(1)\ge1$. By superadditivity,
\[
f(3)\ge3f(1)\ge3.
\]
Also, submultiplicativity gives
\[
f(1)=f\left(3\cdot\frac13\right)
\le f(3)f\left(\frac13\right),
\]
and hence
\[
f\left(\frac13\right)\ge\frac{f(1)}{f(3)}\ge\frac1{f(3)}.
\]
Writing $a=f(3)\ge3$, we obtain
\[
\frac{10}{3}
\ge f(3)+f\left(\frac13\right)
\ge a+\frac1a
\ge 3+\frac13=\frac{10}{3}.
\]
Equality holds throughout, so
\[
f(1)=1,\qquad f(3)=3,\qquad f\left(\frac13\right)=\frac13.
\]

\medskip
\noindent\textbf{Step 2. Establish an upper bound.}

By induction, superadditivity and submultiplicativity give
\[
f(3^n)=3^n\qquad(n=0,1,2,\ldots).
\]
The function is strictly increasing: if $u>v>0$, then
\[
f(u)=f\bigl(v+(u-v)\bigr)
\ge f(v)+f(u-v)>f(v).
\]

Fix $t>0$, and for all sufficiently large $n$ let
\[
m_n=\left\lfloor\frac{3^n}{t}\right\rfloor.
\]
Since $m_nt\le3^n$, monotonicity and superadditivity yield
\[
3^n=f(3^n)\ge f(m_nt)\ge m_nf(t).
\]
As $3^n/m_n\to t$, it follows that
\[
f(t)\le t\qquad(t>0). \tag{1}
\]

\medskip
\noindent\textbf{Method I. Rational squeeze.}

For every positive integer $k$, superadditivity and (1) give
\[
k=kf(1)\le f(k)\le k,
\]
so $f(k)=k$. If $r=p/q$ is a positive rational, where $p,q$ are
positive integers, submultiplicativity gives
\[
p=f(p)=f(qr)\le f(q)f(r)=qf(r).
\]
Thus $f(r)\ge r$, and (1) shows $f(r)=r$. By monotonicity and the
density of the positive rationals, $f(x)=x$ for all $x>0$.

\end{proof}

The identity function satisfies the original conditions with equality
throughout, so it is the unique solution.

\end{dlsummary}